% Generated draft: requires mathematical and source review.
\documentclass[11pt]{article}
\usepackage[margin=25mm]{geometry}
\usepackage{fontspec}
\setmainfont{Latin Modern Roman}
\usepackage{amsmath,amssymb,mathtools,unicode-math}
\setmathfont{Latin Modern Math}
\usepackage{enumitem,needspace,xurl,hyperref}
\setlength{\parindent}{0pt}
\setlength{\parskip}{5pt}
\setlength{\emergencystretch}{3em}
\newcommand{\N}{\mathbb{N}}
\newcommand{\Z}{\mathbb{Z}}
\newcommand{\Q}{\mathbb{Q}}
\newcommand{\R}{\mathbb{R}}
\newcommand{\C}{\mathbb{C}}
\newcommand{\F}{\mathbb{F}}
\newcommand{\A}{\mathbb{A}}
\newcommand{\PP}{\mathbb{P}}
\newcommand{\E}{\mathbb{E}}
\newcommand{\eps}{\varepsilon}
\newcommand{\dd}{\,\mathrm{d}}
\begin{document}
\section*{UnsolvedMath ID: 800058}
\textbf{Source problem code:} catalogue entry
\subsection{Definitions and mathematical statement}
Let \(S^2=\{x\in\R^3:\|x\|=1\}\), and let \(M=S^2\times S^2\). Under the usual interpretation that \(M\) carries its standard product smooth structure, a complex structure on \(M\) is an integrable smooth bundle endomorphism \(J:TM\to TM\) satisfying \(J^2=-\operatorname{id}_{TM}\). Integrability means that \(J\) is induced by an atlas of complex coordinate charts with holomorphic transition maps; the resulting complex manifold has complex dimension two.

\textbf{Classification problem.} Determine all complex structures on \(S^2\times S^2\).

\textbf{Unresolved qualifications.} The supplied statement does not specify an equivalence relation for this classification. If classification up to biholomorphism is intended, the target is
\[
\bigl\{X:\ X\text{ is a compact connected complex surface and its underlying smooth manifold is diffeomorphic to }M\bigr\}/\cong_{\mathrm{bihol}},
\]
where \(X\cong_{\mathrm{bihol}}Y\) means that there is a holomorphic bijection \(X\to Y\) with holomorphic inverse. Equivalently, complex structures \(J_1,J_2\) on the fixed smooth manifold are identified when there is a diffeomorphism \(f:M\to M\) such that \(\dd f\circ J_1=J_2\circ\dd f\). This is a conditional interpretation, not an equivalence relation explicitly supplied by the source. The source also does not explicitly distinguish the standard smooth-product interpretation from a classification of complex surfaces merely homeomorphic to \(M\), nor does it impose compatibility with a fixed orientation.
\subsection{Short English statement}
Determine all complex surface structures on \(S^2\times S^2\). The intended equivalence relation and underlying smooth or topological interpretation require clarification. The supplied registry labels the problem open; this draft does not independently verify that status.
\subsection{Sources}
\begin{itemize}
\item Local UnsolvedMath registry, supplied complete source record: permanent ID \(800058\), source problem code \texttt{catalogue entry}, title ``Hirzebruch's 1954 Problem List --- Problem 7.'' The record attributes the entry to Kotschick, with proposed year 2013, and labels it \texttt{open}.
\item Kotschick, \emph{Updates on Hirzebruch's 1954 Problem List} (2013), Problem 7. Supplied source URL: \url{https://arxiv.org/abs/1305.4623}. The record identifies the source list with Hirzebruch's 1954 problem list.
\item Supplied provenance reports access on 2026-10-16 and extraction by \texttt{arxiv\_latex}. These are supplied provenance details, not a claim of independent source access or verification. The registry record was created and updated on 2026-10-16.
\end{itemize}
\end{document}
